% Einheit 3 von 3 – Sonderansätze und Modellkritik (bank/rekonstruktion-von-funktionsgleichungen/e3.jsonl, zusammenbau v0.1)
%% TODO Zeitmarke Einheit 3: Marken-Zeile nicht lesbar
%% TODO Zweigzeile Teil 1: was hier gelernt wird (Satz in Schülersprache aus dem Einheitstitel) – Einheit 3
\einheitenkopf[e3]{Einheit 3 von 3 · Sonderansätze und Modellkritik}
\zweigzeile{Abitur GK}
\setcounter{aufgabe}{13}

%% TODO Titel: Ich-kann-Satz fehlt in der Bank (Platzhalter: Kettenname) – Nr. 14
%% TODO Anweisung: ein Satz über der ersten Teilaufgabe fehlt in der Bank – Nr. 14
\begin{aufgabe}{Sonderansätze und Modellkritik}
%% TODO \rechenplatz{n}: Schrittzahl des Grundfalls fehlt in der Bank (nur bei fester Schreibform, 2.3 b)
\begin{teile}
\teil Ein Sinusgraph hat einen Hochpunkt bei $x = 1$ und den nächsten Tiefpunkt bei $x = 4$. Wie weit ist es vom Hochpunkt bis zum nächsten Hochpunkt? \\ \kreuz{$3$} \\ \kreuz{$6$} \\ \kreuz{$12$}
\teil Der Graph von $f$ mit $f(x) = a \cdot \mathrm{sin}(b \cdot x)$ und $a, b > 0$ hat die direkt aufeinanderfolgenden Extrempunkte $E_1(1|3)$ und $E_2(3|-3)$. Bestimme $a$ und $b$. a = \leerfeld , b = \leerfeld
\teil Die Funktion $f$ mit $f(x) = x - \frac{1}{27}x^3$ soll auf $[-3; 3]$ durch $g$ mit $g(x) = a \cdot \mathrm{sin}(b \cdot x)$ und $a, b > 0$ angenähert werden. Dabei ist $3$ die erste positive Extremstelle von $g$, und es gilt $g(3) = f(3)$. Bestimme $a$ und $b$. \hfill (Abitur 2022 GK) \\ a = \leerfeld , b = \leerfeld
\teil Die Funktion $g$ mit $g(x) = a \cdot \mathrm{sin}(b \cdot x)$ und $a, b > 0$ soll die Nullstellen $0$ und $4$ haben, dazwischen keine weitere. Zwischen diesen Nullstellen soll ihr Graph mit der $x$-Achse eine Fläche mit dem Inhalt 8 einschließen. Bestimme $a$ und $b$. a ≈ \leerfeld , b ≈ \leerfeld
\teil Die Einwohnerzahl einer Stadt wird ab 2020 durch eine quadratische Funktion $p$ beschrieben ($x$ in Jahren seit 2020, $p(x)$ in Tausend). 2020 hat die Stadt 80 Tausend Einwohner. Zwischen 2020 und 2060 ist das durchschnittliche Wachstum null, und die größte Einwohnerzahl liegt um 12 Tausend über der von 2020. Ermittle, wann die Einwohnerzahl am größten ist und wie groß sie dann ist, und gib eine Gleichung von $p$ an. x = \leerfeld , Maximum \leerfeld , p(x) = \leerfeld
\teil Ein Deich wird im Querschnitt beschrieben; eine Längeneinheit entspricht 5 m. Der Querschnitt beginnt auf der $y$-Achse in 9 m Höhe waagerecht und fällt dann bis zum Boden ($x$-Achse) ab; dieser Teil wird durch eine quadratische Funktion $p$ beschrieben. Begründe, dass $p(x) = ax^2 + c$ ohne lineares Glied angesetzt werden kann. Die Fläche zwischen dem Graphen von $p$, der $y$-Achse und der $x$-Achse beträgt 90 m²; bestimme $a$ und $c$. a = \leerfeld , c = \leerfeld
\teil Das Profil einer Kinderrutsche soll auf $[0; 6]$ bei $x = 0$ ohne Knick an eine waagerechte Plattform anschließen (Hochpunkt $(0|3)$) und bei $x = 6$ waagerecht in den Boden auslaufen. Begründe, dass eine einzige quadratische Parabel dieses Profil auf $[0; 6]$ nicht beschreiben kann. \hfill (Abitur 2023 LK)
\end{teile}
\end{aufgabe}

%% TODO Titel: Ich-kann-Satz fehlt in der Bank (Platzhalter: Kettenname) – Nr. 15
%% TODO Anweisung: ein Satz über der ersten Teilaufgabe fehlt in der Bank – Nr. 15
\begin{aufgabe}{Zwei Parameter einer Funktion aus einer vorgegebenen Bogenlängenformel und einem Punkt bestimmen}
\begin{teile}
\teil Eine Fußgänger-Hängebrücke ist 200 m lang; in einem Koordinatensystem (eine Einheit entspricht 1 m) liegt die Mitte auf der $y$-Achse, das Seil ist in 40 m Höhe an den Enden befestigt. Es soll durch $h(x) = \frac{s}{2}\left(e^{x/s} + e^{-x/s}\right) + t$ mit $s > 0$ beschrieben werden und 210 m lang sein. Die Länge des Graphen von $h$ zwischen $(-v|h(-v))$ und $(v|h(v))$ ist $s \cdot \left(e^{v/s} - e^{-v/s}\right)$ für $v > 0$. Bestimme $s$ und $t$ mit dem Rechner auf eine Stelle. \hfill (Abitur 2018 LK) \\ s ≈ \leerfeld , t ≈ \leerfeld
\end{teile}
\end{aufgabe}

%% TODO Titel: Ich-kann-Satz fehlt in der Bank (Platzhalter: Kettenname) – Nr. 16
%% TODO Anweisung: ein Satz über der ersten Teilaufgabe fehlt in der Bank – Nr. 16
\begin{aufgabe}{Existenz einer quadratischen Funktion zu vier Wert- und Steigungsbedingungen über das überbestimmte Gleichungssystem untersuchen}
\begin{teile}
\teil Der Graph einer Funktion $f$ verläuft durch $A(0|3)$ und $B(1|0{,}5)$ mit $f'(0) = -2$ und $f'(1) = 1$. Der Graph einer quadratischen Funktion $p$ soll in $A$ und in $B$ tangential zum Graphen von $f$ verlaufen. Stelle vier Bedingungen für $p$ auf und untersuche, ob es eine solche Funktion $p$ gibt. \hfill (Abitur 2017 GK)
\end{teile}
\end{aufgabe}

%% TODO Titel: Ich-kann-Satz fehlt in der Bank (Platzhalter: Kettenname) – Nr. 17
%% TODO Anweisung: ein Satz über der ersten Teilaufgabe fehlt in der Bank – Nr. 17
\begin{aufgabe}{Unmöglichkeit waagerechter Tangenten an den Rändern über die Ableitung mit Parametern begründen}
\begin{teile}
\teil Ein Profil wird durch $q(x) = a + c \cdot e^{-\frac{x^2}{2}}$ mit $a, c > 0$ und $x \in [-2; 2]$ beschrieben. Begründe, dass $a$ und $c$ nicht so gewählt werden können, dass der Graph an seinen Randpunkten waagerecht ausläuft. \hfill (Abitur 2017 LK)
\end{teile}
\end{aufgabe}

%% TODO Titel: Ich-kann-Satz fehlt in der Bank (Platzhalter: Kettenname) – Nr. 18
%% TODO Anweisung: ein Satz über der ersten Teilaufgabe fehlt in der Bank – Nr. 18
\begin{aufgabe}{Sonderansätze und Modellkritik – Fehler finden, Begründen}
\begin{teile}
\teil Ich finde den Fehler: Der Graph von $f$ mit $f(x) = a \cdot \mathrm{sin}(b \cdot x) + 1$ und $a, b > 0$ hat die direkt aufeinanderfolgenden Extrempunkte $H(2|4)$ und $T(6|-2)$. Ole rechnet: \rechnung{p &= 6 - 2 = 4 \\ b &= \frac{2\pi}{4}}
\teil Begründe, warum ein Hochpunkt und der nächste Tiefpunkt eines Sinusgraphen eine halbe Periode auseinanderliegen.
\end{teile}
\end{aufgabe}
